\chapter*{Kết luận và hướng phát triển}
\addcontentsline{toc}{chapter}{Kết luận và hướng phát triển}

Khóa luận đã giải quyết hai bài toán có liên hệ mật thiết với nhau:
nghiên cứu phương pháp học máy để phát hiện mã độc PE và triển khai
kết quả nghiên cứu vào một hệ thống bảo vệ endpoint hoạt động được
trên Windows.

Về bài toán nghiên cứu, khóa luận đã xây dựng và đánh giá thực nghiệm
ba pipeline học máy trên tập dữ liệu EMBER2024, tập trung vào việc
đánh giá đóng góp của kỹ thuật trích xuất năng lực mã độc (capability
extraction) bằng công cụ Capa. Kết quả cho thấy tất cả các pipeline
có bổ sung đặc trưng Capa đều vượt baseline EMBER2024 thuần túy trên
Partial AUC (FPR $\leq$ 0,5\%). Pipeline kết hợp EMBER2024 với Capa
đã giảm chiều bằng TruncatedSVD-256 đạt kết quả tốt nhất với Partial
AUC = 0,8850 và Full AUC = 0,9926, cải thiện lần lượt 0,34\% và
0,04\% so với baseline (0,8816 và 0,9922). Kết quả này xác nhận rằng
thông tin ngữ nghĩa hành vi từ Capa mang lại giá trị phân biệt bổ
sung thực chất, đồng thời cho thấy TruncatedSVD vừa phù hợp với đặc
trưng dạng thưa của Capa, vừa có tác dụng lọc nhiễu tốt hơn so với
multi-hot encoding thuần túy.

Về bài toán công nghệ, khóa luận đã thiết kế và hiện thực hóa kiến
trúc hệ thống bảo vệ endpoint gồm năm thành phần: module trích xuất
đặc trưng C++ (ember2024-fe-cpp) đạt tốc độ nhanh hơn cài đặt Python
gốc từ 12 đến 20 lần; kernel driver phát hiện tiến trình thực thi
(pmd-driver); dịch vụ nền tích hợp mô hình học máy (pmd-service); giao
diện người dùng với thông báo thời gian thực (pmd-ui); và dịch vụ phân
tích sâu trực tuyến tích hợp Capa và SHAP (deep-analysis-server). Việc
tích hợp SHAP theo nhóm đặc trưng giúp người dùng hiểu được lý do mô
hình đưa ra quyết định phân loại, nâng cao tính minh bạch và khả năng
ứng dụng thực tế của hệ thống.

Tuy nhiên, khóa luận vẫn còn một số hạn chế. Về
phía nghiên cứu, thực nghiệm so sánh các pipeline được thực hiện trên
tập con 450.000 mẫu thay vì toàn bộ tập EMBER2024, do giới hạn về
thời gian chạy Capa. Đề tài cũng chỉ tập trung vào phân tích tĩnh và
một mô hình học máy duy nhất (LightGBM), chưa so sánh với các kiến
trúc học sâu. Về phía hệ thống, driver chưa được ký số nên chỉ có thể
nạp qua chế độ Provisioning trên máy ảo phát triển (dev machine), không phù hợp với
triển khai trên máy người dùng thực tế (song đây là giới hạn của
phạm vi đề tài). Ngoài ra, kernel module hiện chưa có các cơ chế chống
can thiệp (anti-tampering) cần thiết cho môi trường triển khai thực tế.

Từ những kết quả và hạn chế trên, khóa luận đề xuất một số hướng
phát triển tiếp theo. Trước hết, cần mở rộng thực nghiệm trên toàn
bộ tập EMBER2024 và so sánh với các kiến trúc học sâu để có đánh giá
toàn diện hơn về hiệu quả của capability extraction. Tiếp theo, mặc
dù bổ sung đặc trưng Capa mang lại cải thiện có ý nghĩa, mức cải
thiện còn khiêm tốn gợi ý rằng việc biểu diễn capability ở mức hành
vi đơn lẻ chưa khai thác hết tiềm năng của thông tin ngữ nghĩa. Một
hướng nghiên cứu đầy triển vọng là nâng mức trừu tượng từ capability
lên \textit{intent} - tức là suy luận ý đồ của mã độc (đánh cắp dữ
liệu, duy trì/persistence, tránh phát hiện/evade detection...)
từ sự kết hợp của
nhiều capability, thay vì xử lý từng capability độc lập. Đây là một
research gap còn khá mở, trong đó các framework như MITRE ATT\&CK có
thể đóng vai trò nền tảng để xây dựng biểu diễn đặc trưng cấp cao
hơn. Song song đó, tích hợp phân tích động theo hướng hybrid analysis
- ví dụ thông qua knowledge distillation như đã trình bày trong
Chương~1 - có thể nâng cao khả năng phát hiện các mẫu mã độc trốn
tránh phân tích tĩnh. Về phía hệ thống, các hướng phát triển
tiếp theo bao gồm cho phép cấu hình ngưỡng phân loại (threshold) linh hoạt;
tăng cường cơ chế chống can thiệp tại kernel module
(pmd-driver) nhằm đảm bảo tính toàn vẹn trong
môi trường thực tế, bao gồm tận dụng Trusted Platform Module (TPM);
bổ sung callback ở cấp kernel để giám sát thêm các sự kiện như DLL
image mapping, và xây dựng cơ chế tạm thời khóa file PE trong quá
trình chờ quét ở user mode. Cuối cùng, mở rộng pipeline sang các định
dạng file khác ngoài PE - chẳng hạn ELF, PDF và APK - là các
hướng phát triển tiềm năng khi nền tảng nghiên cứu đã được
xây dựng vững chắc.

Nhìn chung, khóa luận đã có những đóng góp ở cả phương diện
nghiên cứu và ứng dụng thực tiễn. Về mặt nghiên cứu, đề tài
cung cấp các bằng chứng thực nghiệm cho thấy hiệu quả của
phương pháp trích xuất capability trong bài toán phát hiện
mã độc PE. Về mặt công nghệ, khóa luận xây dựng thành công
một hệ thống bảo vệ endpoint hoàn chỉnh, kết hợp giữa lập
trình hệ thống trên Windows và các kỹ thuật học máy hiện đại.
Những kết quả và kinh nghiệm thu được từ đề tài được kỳ vọng
sẽ trở thành nền tảng hữu ích cho các nghiên cứu và ứng dụng
tiếp theo trong lĩnh vực phát hiện mã độc và an toàn thông tin.
